\chapter{Kesimpulan dan Saran}

\section{Kesimpulan}

Penelitian ini telah berhasil merancang dan mengimplementasikan peripheral I2S TX-only
berbasis Verilog dengan antarmuka AXI4-Lite pada papan pengembangan Basys3
(Xilinx Artix-7 XC7A35T). Berdasarkan hasil implementasi dan pengujian yang dilakukan,
dapat disimpulkan bahwa:

\begin{enumerate}
\item Peripheral I2S TX-only berhasil dirancang dalam bahasa Verilog dengan
      antarmuka AXI4-Lite yang mencakup sembilan register yang dapat dikonfigurasi
      melalui MicroBlaze, menghasilkan sinyal BCLK, WS, dan DATA yang memenuhi
      standar I2S Philips dengan slot 32-bit dan data audio 24-bit MSB-\textit{first}.

\item Dua mode laju sampel berhasil diimplementasikan menggunakan satu bit kendali
      FS\_MODE pada register CONTROL:
      \begin{itemize}
      \item Mode~1 ($\approx$96\,kHz): BCLK $\approx$6,142\,MHz, WS $\approx$95,97\,kHz.
      \item Mode~0 ($\approx$48\,kHz): BCLK $\approx$3,071\,MHz, WS $\approx$47,99\,kHz.
      \end{itemize}
      Galat frekuensi aktual terhadap nilai target ($\pm$0,029\%) berada dalam batas
      toleransi yang dapat diterima DAC PCM5102A.

\item TX FIFO berkedalaman 256 frame stereo berhasil diimplementasikan sebagai array
      register sinkron, dengan mekanisme \textit{push} pada penulisan TX\_RIGHT dan
      \textit{pop} pada setiap batas frame I2S. Kondisi \textit{underrun} ditangani
      dengan keluaran nol otomatis dan indikasi flag status.

\item Mekanisme IRQ \textit{watermark} yang dapat dikonfigurasi berhasil diimplementasikan
      dan diuji dengan aplikasi demo \texttt{main.cpp} pada Xilinx Vitis, memungkinkan
      ISR mengisi ulang FIFO secara efisien tanpa \textit{polling} konstan.

\item Fungsionalitas sistem berhasil diverifikasi melalui dua metode: simulasi RTL
      menggunakan Vivado Simulator dan pengujian perangkat keras langsung menggunakan
      ILA pada Basys3 dengan keluaran audio nyata ke DAC PCM5102A. Sinyal sinus
      1\,kHz terputar bersih tanpa artefak yang terdengar.
\end{enumerate}

\section{Saran}

Berdasarkan keterbatasan penelitian yang ada, berikut adalah saran untuk pengembangan
lebih lanjut:

\begin{enumerate}
\item \textbf{Integrasi DMA (\textit{Direct Memory Access})}: Menggantikan mekanisme
      pengisian FIFO berbasis IRQ dari CPU (Vitis \texttt{main.cpp}) dengan AXI DMA
      dari Xilinx. Data audio dapat disimpan di BRAM atau DDR, kemudian dipindahkan
      ke FIFO secara otomatis tanpa intervensi CPU, sehingga membebaskan CPU untuk
      tugas lain.

\item \textbf{Dukungan Format Lebih Lanjut}: Menambahkan dukungan untuk format
      \textit{left-justified} (LJ) dan \textit{right-justified} (RJ) selain format
      standar Philips, serta mode TDM (\textit{Time Division Multiplexing}) untuk
      lebih dari dua kanal audio.

\item \textbf{Mode Laju Sampel Tambahan}: Menambahkan mode laju sampel 192\,kHz
      (BCLK = audio\_clk/2) dan 44,1\,kHz dengan konfigurasi MMCM yang sesuai, atau
      menggunakan osilator audio eksternal (24,576\,MHz atau 22,5792\,MHz) untuk
      presisi frekuensi yang lebih baik.

\item \textbf{Penerimaan I2S (RX Mode)}: Menambahkan modul I2S receiver untuk
      mendukung operasi penuh (\textit{full duplex}) dengan mikrofon digital atau
      ADC yang menggunakan protokol I2S, sehingga sistem dapat digunakan untuk
      aplikasi rekaman dan pemrosesan audio real-time.

\item \textbf{Integrasi Linux ASoC Driver}: Pada tahap pengembangan lanjutan
      dengan platform yang mendukung Linux (misalnya Zynq-7000 atau Zynq UltraScale+),
      peripheral ini dapat dijadikan dasar untuk pengembangan \textit{CPU DAI driver}
      menggunakan framework ASoC (\textit{ALSA System on Chip}) dengan
      \textit{Device Tree} binding yang sesuai.

\item \textbf{Pengujian Kualitas Audio}: Melakukan pengukuran kualitas audio yang
      lebih komprehensif seperti SNR (\textit{Signal-to-Noise Ratio}), THD
      (\textit{Total Harmonic Distortion}), dan SINAD menggunakan perangkat lunak
      analisis audio (misalnya REW atau ARTA) untuk memvalidasi kualitas keluaran
      secara objektif.
\end{enumerate}
